\repchapternew{18}{Загальний аналітичний висновок}

У цьому розділі подано авторське аналітичне судження, сформульоване на підставі матеріалів серії Social з урахуванням їхніх застережень. Це висновок за зібраними матеріалами, а не самостійне розслідування всіх первинних джерел. Рівні доказовості позначено за шкалою проєкту: \level{a} --- виміряно, \level{b} --- кілька незалежних джерел, \level{c} --- одне джерело або сильні альтернативи, \level{d} --- без перевірюваних даних. Походження чисел наведено в підрозділі~18.13.

\begin{keybox}{Основна теза}
Подія 2 листопада 2026 року була реальним локальним соціальним конфліктом, що переріс в інформаційну кризу загальноукраїнського масштабу. Спостерігається не одна інформаційна кампанія з доведеним центром управління, а взаємодія реального конфлікту, інформаційної невизначеності, емоційної мобілізації та вибіркового підсилення різними учасниками.
\end{keybox}

\repsection{18.1}{Три одночасні процеси}
Розвиток подій описують три процеси, які не можна звести один до одного.
\begin{enumerate}
\item \textbf{Реальний конфлікт} навколо військовослужбовців, умов їхнього перебування та дій людей біля об'єкта НГУ. Він породив кілька самостійних правових і суспільно значущих питань.
\item \textbf{Інформаційна хвиля}, у якій непідтверджені й емоційно сильні твердження поширювалися швидше, ніж з'являлася достатня перевірена інформація.
\item \textbf{Боротьба інтерпретацій} за участю українських пабліків, офіційних установ, редакційних медіа, каналів із попередніми російськими або проросійськими мітками та звичайних користувачів.
\end{enumerate}
Російське інформаційне середовище, наскільки дозволяють судити матеріали, використовувало окремі сюжети для підсилення критичних інтерпретацій українських інституцій. Проте доказів того, що російська держава створила сам інцидент, ініціювала першу хвилю або керувала її учасниками, у зібраних матеріалах немає.

\repsection{18.2}{Реконструкція конфліктної ситуації}
Об'єкт у Петропавлівській Борщагівці використовувався Національною гвардією як тимчасовий пункт розміщення військовослужбовців. За заявою командувача НГУ на брифінгу 3 листопада, там перебували 26 осіб: четверо військових, які самовільно залишили частину й повернулися, та інші, які готувалися до базової загальновійськової підготовки. В інших повідомленнях фігурували 17 і 18 осіб. Наявні документи не дозволяють установити точну кількість людей для кожного моменту часу. На цьому тлі виник конфлікт щодо перебування конкретного військовослужбовця, стану його здоров'я та можливості його забрати.

Існують дві суттєво різні версії. За версією слідства, окремі особи заздалегідь домовилися організувати вивезення військовослужбовця, а публічна активність навколо об'єкта була частиною цього плану; у публікаціях згадуються можливі гроші, попередня розвідка об'єкта та використання медичного транспорту. За версією підозрюваних і захисту, йшлося про спробу допомогти військовослужбовцю, який нібито потребував медичної допомоги.

\emph{Оцінка.} Підстави розглядати можливу попередню організацію дій окремих осіб є, оскільки таку версію офіційно висунуло слідство (рівень \level{c}). Однак відкритих доказів недостатньо, щоб вважати встановленими оплату, попередній злочинний задум або організацію всього натовпу. Це обмеження визначає, який сценарій розвитку конфлікту взагалі допустимо реконструювати. Вину будь-кого з підозрюваних може встановити лише суд.

Первинні повідомлення говорили приблизно про 40 людей у підвалі, катування та загибель людини. Поліція заперечила ці твердження, а незалежних підтверджень такого масштабу подій матеріали не містять. Звідси не випливає, що всі скарги на поводження з військовослужбовцями були безпідставними. Окрема публікація пресслужби НГУ за матеріалами АрміяInform повідомила про чотирьох військовослужбовців із травмами, яких направили до госпіталю; про чотирьох травмованих повідомляли також корпус «Азов» і МВС. Час отримання травм, їхні обставини й причини лишаються невстановленими, а доступних медичних документів недостатньо для незалежного висновку. Крім того, ДБР відкрило окреме провадження щодо можливого перевищення влади військовою службовою особою, МВС оголосило службове розслідування, а офіс військової омбудсманки призначив перевірку.

Отже, існують два окремі питання, які не можна підміняти одне одним: чи були неправомірними дії людей, що зібралися біля об'єкта НГУ, і чи порушувалися права військовослужбовців усередині об'єкта. Перевірка одного не скасовує необхідності перевірити друге. Змішування цих питань, на думку автора, стало однією з причин гостроти інформаційного конфлікту.

\repsection{18.3}{Розгортання інформаційної хвилі}
У таблиці~\ref{tab:chron} наведено хронологію публікацій у доступних джерелах. Це не доведений причинний ланцюг: перший знайдений допис у Telegram не обов'язково є першою публікацією взагалі, а можливі раніші повідомлення в Threads відомі лише ретроспективно.

\begin{longtable}{@{}>{\raggedright\arraybackslash}p{0.13\linewidth}>{\raggedright\arraybackslash}p{0.44\linewidth}>{\raggedright\arraybackslash}p{0.37\linewidth}@{}}
\caption{Хронологія 2 листопада 2026 року та її аналітичне значення (час київський).}\label{tab:chron}\\
\toprule\textbf{Час} & \textbf{Зафіксована подія} & \textbf{Аналітичне значення}\\\midrule\endfirsthead
\toprule\textbf{Час} & \textbf{Зафіксована подія} & \textbf{Аналітичне значення}\\\midrule\endhead
\bottomrule\endfoot
\footnotesize 11:34 & \footnotesize За ретроспективним повідомленням ЦПД, у Threads з'явився прямий заклик приїхати до будівлі; первинний допис не отримано & \footnotesize Можлива рання точка онлайн-мобілізації, але не встановлене першоджерело\\
\footnotesize 12:41:03 & \footnotesize Перший релевантний допис у панелі Telegram: приблизно 40 людей у підвалі «зі слів медиків» & \footnotesize Поява емоційно сильного повідомлення з посиланням на нібито свідчення медиків\\
\footnotesize 13:07:53 & \footnotesize Публікація з формулюванням «Будні концтабору» & \footnotesize Перехід від конкретного твердження до різко негативної інтерпретації\\
\footnotesize 13:16:21 & \footnotesize Повідомлення публікує великий канал «Реальний Київ» & \footnotesize Розширення потенційної аудиторії сюжету\\
\footnotesize 13:57:01 & \footnotesize Публікація «Украина.ру» з посиланням на київські ЗМІ & \footnotesize Підхоплення сюжету російським медіа після його появи в українських джерелах панелі\\
\footnotesize 17:06--22:47 & \footnotesize У Telegram поширюють посилання на прямі ефіри TikTok & \footnotesize Відео й трансляції стають важливою частиною інформаційного середовища\\
\footnotesize 21:32 & \footnotesize Час спростування, використаний у дослідженні твердження про «медика» & \footnotesize Зміна публічного інформаційного контексту\\
\footnotesize 22:30 & \footnotesize Найвищий десятихвилинний пік публікацій у панелі Telegram (32 дописи) & \footnotesize Обговорення лишалося інтенсивним і після появи спростувань\\
\end{longtable}

На першому етапі формується сильна інформаційна рамка. Її складають шість елементів: \emph{жертви} (люди, яких нібито незаконно утримують), \emph{масштаб} (приблизно 40 осіб), \emph{моральний тригер} (повідомлення про катування та смерть), \emph{джерело авторитету} (посилання на медиків), \emph{інституційний антагоніст} (військові органи й ширше --- система мобілізації) і \emph{заклик до дії} (потреба приїхати й допомогти). У такій конструкції читач отримує не лише новину, а ситуацію, де бездіяльність може здаватися морально неприйнятною. Це не доводить, що автори свідомо її спроєктували, але вона об'єктивно має властивості, здатні сприяти швидкій мобілізації уваги.

\begin{keybox}{Перший аналітичний висновок}
Найсильнішим механізмом ранньої хвилі була не складна технологічна маніпуляція, а поєднання емоційно гострого повідомлення, претензії на свідчення очевидців і суспільно чутливої теми.
\end{keybox}

\repsection{18.4}{Чинники швидкого масштабування}
\emph{Довіра.} Для суспільства у воєнний час питання мобілізації, поводження з військовослужбовцями, повноважень державних органів і прав людини є особливо чутливими. Це контекстуальний висновок, а не результат вимірювання: репрезентативного опитування матеріали не містять. Проте зміст повідомлень добре узгоджується з тим, що конфлікт потрапив у вже наявне поле суспільної напруги.

\emph{Візуальні свідчення.} Відео, медичний одяг, автомобіль із медичним брендуванням і прямі ефіри могли створювати враження безпосереднього доступу до події, однак візуальна переконливість не дорівнює достовірності. Питання особистої кваліфікації В.~Молодої, правомірності використання медичного транспорту, ліцензій і статусу компаній складніші за ярлики «справжня» чи «фальшива» медикиня. Ліцензія компанії не встановлює повноважень конкретної особи працювати бригадою екстреної допомоги; відсутність ліцензії на окремий вид практики не доводить відсутності медичної освіти.

\emph{Зв'язок онлайн-повідомлень із рішеннями учасників.} За судовим репортажем, В.~Ладика заявила, що приїхала на місце після перегляду ефірів TikTok. Це самозвіт однієї учасниці про власний мотив. Він не доводить, що так прибули всі чи більшість присутніх, але показує механізм, за якого інформаційне повідомлення може стати безпосередньою підставою для рішення приїхати (рівень \level{c}). Цей епізод пов'язує інформаційну та фізичну складові конфлікту, не дозволяючи водночас оцінити масштаб причинного впливу.

\repsection{18.5}{Конкуренція інтерпретацій}
Сторони боролися не лише за встановлення фактів, а й за право визначати, що означає сама подія. Умовно сформувалися три способи тлумачення (табл.~\ref{tab:frames}).

\begin{longtable}{@{}>{\raggedright\arraybackslash}p{0.26\linewidth}>{\raggedright\arraybackslash}p{0.33\linewidth}>{\raggedright\arraybackslash}p{0.35\linewidth}@{}}
\caption{Три інтерпретаційні рамки та межі кожної.}\label{tab:frames}\\
\toprule\textbf{Рамка} & \textbf{Фактична основа} & \textbf{Межа}\\\midrule\endfirsthead
\toprule\textbf{Рамка} & \textbf{Фактична основа} & \textbf{Межа}\\\midrule\endhead
\bottomrule\endfoot
\footnotesize «Держава незаконно утримує та катує людей» & \footnotesize Скарги на поводження з військовослужбовцями, повідомлення про можливі порушення й травми, недовіра до офіційних пояснень & \footnotesize Найдраматичніші твердження (40 незаконно утримуваних, масові катування, загибель) належного підтвердження не отримали; «концтабір ТЦК» є оціночною рамкою\\
\footnotesize «Була заздалегідь організована провокація» & \footnotesize Версія слідства про попередні домовленості та спробу вивезення військовослужбовця & \footnotesize Потребує перевірки; навіть підтвердження щодо окремих осіб не означало б організованості всіх присутніх\\
\footnotesize «Ситуацію використали для інформаційної дестабілізації» & \footnotesize Певні канали справді поширювали драматичні сюжети й критику державних інституцій & \footnotesize Поширення вигідного комусь сюжету не тотожне виконанню отриманого завдання\\
\end{longtable}

Ці рамки не обов'язково виключають одна одну. Можуть одночасно існувати реальні скарги й невстановлені проблеми, неправдиві або перебільшені твердження, протиправні дії окремих учасників і використання резонансу сторонніми акторами. Саме така багаторівнева модель найкраще пояснює доступну картину.

\repsection{18.6}{Роль російських і проросійських ресурсів}
\emph{Початок хвилі.} У зібраній панелі Telegram перший релевантний допис з'явився о 12:41; перша публікація російського медіа «Украина.ру» зафіксована о 13:57 і сама посилалася на київські ЗМІ (рівень \level{a} для панелі). Це узгоджується з моделлю, у якій наявний сюжет підхоплює інший інформаційний сегмент. Поширити цей висновок на весь інтернет не можна, бо повної хронології Threads, Facebook, TikTok та інших платформ немає. Найкраще підтверджена роль російського медіа в цьому епізоді --- подальше поширення сюжету, а не доведене його створення.

\emph{Концентрація.} За попередньою класифікацією каналів 95,7\% суми знімкових лічильників переглядів дописів із міткою N3 «концтабір» припадали на групу C; після вилучення спірних та автоматичних міток --- 88,9\%. Цей показник має чотири обмеження: група C містила багато каналів, класифікованих автоматично за непрямими ознаками; мітка N3 фіксує згадку сюжету, не обов'язково підтримку; лічильники переглядів не є кількістю унікальних людей; два канали забезпечували 40,8\% відповідних лічильників групи. Тому коректне формулювання таке: у цільовій панелі Telegram поширення повідомлень, пов'язаних із рамкою «концтабору», було сильно сконцентроване серед каналів із попередньою міткою C і особливо залежало від кількох великих джерел. Твердження про «95,7\% російського впливу» дані не підтримують.

\emph{Вибірковість.} У перевірених каналах із ручною попередньою міткою C не знайдено поширення спростування історії про «медика»; натомість траплялися контрспростування, висміювання офіційної реакції та твердження про інформаційне прикриття. Частина цих каналів підтримувала рамку, у якій офіційним поясненням не варто довіряти. Це може бути редакційною позицією, ідеологічною орієнтацією, конкуренцією за увагу або складовою свідомого інформаційного впливу; для встановлення останнього потрібні додаткові докази.

\emph{Питання про операцію.} На рівні всього інциденту російську інформаційну операцію не встановлено. Твердження ЦПД про приблизно 110 російських каналів з однаковим заголовком не вдалося відтворити в наявній панелі. Це не спростовує можливості такої активності поза панеллю, але не дозволяє незалежно використати це твердження як підтвердження спільного керування. Достатніх доказів спільного оператора, завдань, фінансування чи прямого державного контролю матеріали не містять (табл.~\ref{tab:attr}).

\begin{longtable}{@{}>{\raggedright\arraybackslash}p{0.47\linewidth}>{\raggedright\arraybackslash}p{0.38\linewidth}>{\centering\arraybackslash}p{0.09\linewidth}@{}}
\caption{Оцінка окремих питань атрибуції.}\label{tab:attr}\\
\toprule\textbf{Питання} & \textbf{Висновок} & \textbf{Рівень}\\\midrule\endfirsthead
\toprule\textbf{Питання} & \textbf{Висновок} & \textbf{Рівень}\\\midrule\endhead
\bottomrule\endfoot
\footnotesize Чи поширювали російські медіа сюжет? & \footnotesize Так, підтверджено для конкретних джерел & \footnotesize \level{a}\\
\footnotesize Чи використовували окремі канали критичні інтерпретації? & \footnotesize Так, спостерігається в корпусі & \footnotesize \level{a}\\
\footnotesize Чи був інформаційний резонанс вигідним Росії? & \footnotesize Обґрунтована стратегічна інтерпретація, не доказ операції & \footnotesize \level{c}\\
\footnotesize Чи діяла єдина російська координована мережа? & \footnotesize Недостатньо доказів & \footnotesize \level{d}\\
\footnotesize Чи організувала Росія сам конфлікт? & \footnotesize Не встановлено & \footnotesize \level{d}\\
\footnotesize Чи було державне російське керування? & \footnotesize Не встановлено & \footnotesize \level{d}\\
\end{longtable}

\repsection{18.7}{Шаблонні коментарі на підтримку посадовця}
Дослідження виявило 438 шаблонних коментарів від 367 акаунтів під чотирма публікаціями із заявами К.~Буданова. Вони містили повторювані формулювання похвали й мали характерний розподіл між обговореннями. Перестановочні тести показали, що окремі особливості кластера малоймовірні за застосованих нульових моделей: скориговане за Голмом значення становить $p=0{,}0016$ (рівень \level{a}). Це значення є найменшим досяжним для обраних 9\,999 перестановок і 16 тестів, тож воно вказує на відхилення, а не на його величину. Воно також не встановлює, хто організував активність і чи була оплата.

Значення цього результату полягає в тому, що він руйнує спрощене протиставлення «російські канали маніпулюють, українська аудиторія реагує органічно». Ознаки незвичної узгодженості можуть виникати й у повідомленнях на підтримку українських посадовців. Це не означає причетності держави; це означає, що інформаційне середовище потрібно досліджувати симетрично, незалежно від змісту повідомлень.

\repsection{18.8}{Офіційна комунікація}
Інформаційна реакція державних інституцій, на думку автора, мала помітні проблеми з темпом, узгодженістю та поясненням різних обставин. Підстав говорити про доведене умисне приховування немає.

\emph{Темп.} Між першим знайденим дописом у Telegram (12:41) та одним із ранніх повідомлень НГУ, переданим Укрінформом о 17:07, минуло приблизно 4 години 26 хвилин. Повнота пошуку не встановлена, тож це не доводить відсутності будь-якої офіційної комунікації до 17:07. Проте значна частина раннього обговорення, ймовірно, розвивалася до появи цього пояснення.

\emph{Узгодженість.} Раннє повідомлення НГУ, передане Укрінформом о 17:07, стверджувало, що травмованих немає; пресофіцер НГУ заперечив і повідомлення про перелом у первинній історії. Наступного дня корпус «Азов» повідомив про чотирьох поранених, а пізніша публікація НГУ --- про чотирьох військовослужбовців із травмами. Ці повідомлення не обов'язково суперечать одне одному, бо можуть описувати різні моменти часу чи різний обсяг відомостей. Але для аудиторії, яка вже бачила твердження про катування, така зміна могла виглядати непослідовною.

\emph{Обсяг спростування.} Доцільно було з самого початку розвести два повідомлення: «твердження про 40 заручників і загиблого не підтверджені» та «скарги щодо умов перебування й можливих травм перевіряються». Тоді спростування першого не сприймалося б як автоматичне заперечення другого. Це рекомендація з погляду кризових комунікацій, а не встановлений факт сприйняття кожної заяви аудиторією. Популярне твердження, що 94\% переглядів накопичилися до спростування, лишається невстановленим, бо Social не має історичних лічильників переглядів.

\repsection{18.9}{Реакція аудиторії}
Корпус містить десятки тисяч коментарів Telegram, а також коментарі YouTube і TikTok. Однак валідованого повного розподілу підтримки конкретних сторін, довіри до офіційних повідомлень, співчуття військовослужбовцям, підтримки чи засудження протестувальників та емоційних реакцій за темами немає, тож твердження про позицію «більшості» були б необґрунтованими.

Спостерігаються різні реакції: обурення повідомленнями про можливе незаконне утримання, сумнів у достовірності першої історії, сприйняття подій як постановки чи провокації. В обговореннях TikTok траплялися як твердження про приховування доказів, так і виправлення географічної плутанини між Петропавлівською та Софіївською Борщагівкою. Аудиторія не була однорідною й не обов'язково приймала одну версію подій, але частки цих позицій ще не виміряні достатньо надійно. Наступним рівнем Social має стати окремий модуль аналізу реакцій аудиторії (Audience Response Intelligence Engine).

\needspace{18\baselineskip}
\repsection{18.10}{Порівняльна оцінка сценаріїв}
\begin{longtable}{@{}>{\raggedright\arraybackslash}p{0.2\linewidth}>{\raggedright\arraybackslash}p{0.74\linewidth}@{}}
\caption{Сценарії розвитку подій за узгодженістю з наявними доказами.}\label{tab:scen}\\
\toprule\textbf{Статус} & \textbf{Зміст сценарію}\\\midrule\endfirsthead
\toprule\textbf{Статус} & \textbf{Зміст сценарію}\\\midrule\endhead
\bottomrule\endfoot
\footnotesize Найкраще узгоджується з доказами & \footnotesize Реальна конфліктна ситуація навколо військовослужбовців і їхнього перебування в пункті НГУ стала основою для швидкого поширення драматичних повідомлень, частина яких містила непідтверджені твердження про масштаби й характер порушень. Повідомлення йшли через локальні канали, великі українські паблики, прямі ефіри та медіа. Пізніше окремі російські ресурси підхопили сюжет у межах власних рамок. Паралельно з'являлися офіційні пояснення, спростування й контрспростування, а частина активності, зокрема шаблонні коментарі, мала статистично незвичні ознаки узгодженості. Точний причинний вплив онлайн-публікацій на кількість і дії присутніх не встановлений. Усі ці процеси могли співіснувати без єдиного центру управління всією хвилею.\\
\footnotesize Лишається відкритим & \footnotesize Окремі учасники заздалегідь планували дії біля об'єкта та використання інформаційного резонансу. Це узгоджується з версією слідства, але потребує перевірки первинних матеріалів. Навіть підтвердження такої організації автоматично не встановить іноземного замовника.\\
\footnotesize Бракує доказів & \footnotesize Російська держава організувала подію, керувала місцевими учасниками та координувала весь інформаційний процес. За наявними матеріалами підтвердити цю версію не можна.\\
\end{longtable}

\repsection{18.11}{Структурний ризик}
Найбільше занепокоєння викликає не сам факт поширення повідомлень російськими каналами і не кількість публікацій, а швидкість, із якою невизначена інформація про реальну подію перетворилася на категоричні, взаємно несумісні твердження. Питання про те, що відбувається з людьми в будівлі, змінилося твердженнями про незаконне утримання, катування й загиблих; потім виникла протилежна категоричність, за якою все є провокацією, а очевидці обманюють; паралельно з'являлися нові обвинувачення, контрспростування й припущення про приховування доказів. У цій послідовності найважливіше питання --- реальні умови перебування військовослужбовців і обставини конфлікту --- ризикує загубитися між конкуруючими версіями.

Це не виміряний психологічний ефект, а аналітична оцінка структури дискусії. Саме така ситуація створює сприятливі умови для інформаційного впливу: достатньо вибірково підсилювати наявні твердження, не вигадуючи всього конфлікту.

\repsection{18.12}{Підсумкове формулювання}
\begin{keybox}{Висновок}
Інцидент у Петропавлівській Борщагівці 2 листопада 2026 року був реальним соціальним конфліктом навколо військового об'єкта, який супроводжувався швидким поширенням непідтверджених тверджень, емоційно забарвлених інтерпретацій та взаємно суперечливих публічних повідомлень.

Доступні матеріали підтверджують факт сутичок, присутність військовослужбовців, поширення драматичних повідомлень про умови їхнього перебування, подальші правоохоронні дії та початок окремих перевірок можливих порушень прав військовослужбовців.

Водночас не встановлено достовірність первинних тверджень про 40 незаконно утримуваних осіб, масові катування й загибель людини. Повідомлення про травмованих військовослужбовців мають офіційну документальну основу у вигляді публікацій НГУ, але час, причини й обставини травм лишаються невстановленими.

Інформаційна хвиля в досліджуваній панелі Telegram почалася з публікацій українських джерел. «Украина.ру» долучилася пізніше, посилаючись на київські ЗМІ. Подальше поширення окремих сюжетів серед каналів із попередніми російськими чи проросійськими мітками свідчить про їхнє інформаційне підсилення в межах вибірки, але не доводить російського ініціювання події чи централізованого державного керування.

Аналіз також виявив статистично незвичну узгодженість окремого кластера коментарів на підтримку українського посадовця, що вказує на необхідність оцінювати можливу координовану поведінку незалежно від політичної спрямованості повідомлень.

Найкраще підтверджена загальна модель --- реальний локальний конфлікт, інформаційне поширення якого відбувалося в умовах неповної перевірки фактів, суспільно чутливих обставин і конкуренції інтерпретацій. Для встановлення організаторів конкретних дій, фактичних обставин можливого порушення прав військовослужбовців та наявності зовнішнього керування потрібні додаткові первинні докази.
\end{keybox}

Кейс виявляє три окремі проблеми. \emph{Правоохоронна й правозахисна:} необхідно встановити, чи були порушення з боку людей біля об'єкта та чи порушувалися права військовослужбовців усередині нього; це незалежні питання, а відповідальність має встановлюватися індивідуально. \emph{Інформаційна:} ранні непідтверджені твердження набули значного поширення, частина каналів вибірково підсилювала певні інтерпретації, а офіційні пояснення не усунули всіх важливих суперечностей. \emph{Інституційна:} виявляється ризик втрати довіри до офіційних пояснень, коли держава не може швидко надати повну, послідовну й перевірювану картину складної події; саму втрату довіри та її масштаб ще потрібно виміряти.

\begin{keybox}{Головний висновок в одному реченні}
Петропавлівська Борщагівка є прикладом того, як реальна конфліктна подія з невстановленими обставинами може стати основою масштабної інформаційної кризи, у якій різні актори підсилюють вигідні їм інтерпретації, а перевірені факти не встигають сформувати спільну картину подій.
\end{keybox}

Формули «російська ІПСО» чи «українська провокація» на поточному рівні доказів були б надмірним спрощенням. Аналітична цінність зібраних матеріалів полягає в можливості показати, як реально розвивався конфлікт, де з'явилися неперевірені твердження, хто їх поширював, які факти з'ясувалися пізніше і що лишилося невідомим.

\repsection{18.13}{Походження чисел}
\begin{longtable}{@{}>{\raggedright\arraybackslash}p{0.33\linewidth}>{\raggedright\arraybackslash}p{0.61\linewidth}@{}}
\caption{Файли результатів і скрипти для чисел цього розділу.}\\
\toprule\textbf{Показник} & \textbf{Файл результатів; скрипт}\\\midrule\endfirsthead
\toprule\textbf{Показник} & \textbf{Файл результатів; скрипт}\\\midrule\endhead
\bottomrule\endfoot
\footnotesize Час перших дописів, пік 22:30 & \footnotesize \path{reports/continuation_20261105/results.json}; \path{src/analysis/analyze.py}\\
\footnotesize 95,7\%, 88,9\% & \footnotesize \path{reports/continuation_20261105/audit_results.json}; \path{src/analysis/continuation_audit.py}\\
\footnotesize 40,8\% & \footnotesize \path{reports/borshchahivka-report.md}, рядок KB12\\
\footnotesize 438 коментарів, 367 акаунтів & \footnotesize \path{reports/telegram-comments/coordinated-comments.md}\\
\footnotesize $p=0{,}0016$ & \footnotesize \path{reports/coordination-null-models-2026-11-06/results.json}; \path{src/analysis/coordination_null_models.py}\\
\footnotesize 17:06--22:47, 17:07, 21:32, 94\% & \footnotesize \path{reports/borshchahivka-report.md}\\
\footnotesize 26, 17 і 18 осіб, чотири травмовані & \footnotesize \path{reports/figures-and-proceedings-2026-11-06.md}\\
\end{longtable}
